% GENERATED FILE. Edit canonical lessons, then run npm run generate:books.
% canonical-source-hash: fnv1a64:ddd2a9dee881f52c
% canonical-lessons: UR-C27-shahr, UR-C27-mulk, UR-C27-lahore, UR-C27-lahori, UR-R27-where-you-are-from

\chapter{Where You Are From}
\label{ch:ur-where-you-are-from}
\hlchaptermodality{27}

\begin{chapteropening}
\textbf{By the end of this chapter:} \emph{I can name a city, a country and a place, and I can say which place I am from.}

The book writes its first place name. With \ur{شہر}, \ur{ملک}, \ur{لاہور} and the nisba -\ur{ی}, a learner can finally answer the two questions every A1 speaking test asks after a name -- and the three words come from Persian, Arabic and Sanskrit respectively, which is the layering of Urdu shown on a map.
\end{chapteropening}

\section[shahr]{\ur{شہر} (shahr) — city}

\label{lesson:UR-C27-shahr}



\begin{quote}
Say \textbf{in the room}. Then say \textbf{here}.

\begin{itemize}
\raggedright
  \item \emph{Recall:} two chapters back, the ending that means plural-and-a-postposition-is-coming — \emph{-oṅ}
\end{itemize}
\end{quote}

\subsection*{You'll want to know: \ur{شہر}}
\begin{quote}
\textbf{\ur{اس} \ur{شہر} \ur{میں}} — \emph{is shahr meṅ} — \textbf{in this city.}
\end{quote}

\begin{quote}
\textbf{\ur{پرانا} \ur{شہر}} — \emph{purānā shahr} — \textbf{the old city.}
\end{quote}

Twenty-six chapters, and the largest place this book has named is a room. \textbf{\ur{شہر}} is the first word in it for somewhere you could be \emph{from}.

\textbf{\ur{شہر}} does not end in \textbf{-\ur{ا}}, so it does not change shape before the postposition — but \textbf{\ur{اس}} does, exactly as it did with \textbf{\ur{کمرے}}.

\begin{cousinweb}
\textbf{\ur{شہر}} is Persian, and behind it stands Old Persian \emph{xsathra}, meaning \textbf{realm}. It is the same word that ends the royal name \emph{Ardashir}, and the same root as Sanskrit \emph{kshatra}, \textquotedblleft{}dominion\textquotedblright{} — which is where \emph{Kshatriya}, the warrior caste, comes from.

So a \textbf{city} in Urdu is etymologically \emph{the place somebody rules}, where \textbf{\ur{کمرہ}} came in by sea from Portugal. Two words for two kinds of space, from two directions, a thousand years apart.
\end{cousinweb}

\subsection*{Your turn}
Say these aloud:

\begin{itemize}
\raggedright
  \item \emph{shahr}
  \item \emph{is shahr meṅ}, then \emph{us shahr meṅ}
  \item \emph{purānā shahr}
  \item the two scales — \emph{is kamre meṅ}, \emph{is shahr meṅ}
\end{itemize}

\begin{tcolorbox}[breakable,colback=teal!4,colframe=teal!35!black,title={Before you move on}]
What did \textbf{\ur{شہر}} originally mean? (\textbf{A realm} — the place somebody rules.) Does it change shape before a postposition? (\textbf{No} — it does not end in \textbf{-\ur{ا}}.) And what does change? (\textbf{The demonstrative in front of it}.)

Source: \href{https://www.unicode.org/charts/PDF/U0600.pdf}{Unicode Arabic chart}.
\end{tcolorbox}
\section[mulk]{\ur{ملک} (mulk) — country}

\label{lesson:UR-C27-mulk}



\begin{quote}
Say \textbf{in this city}. Then say what happened to \textbf{\ur{اس}} to let you say it.
\end{quote}

\subsection*{You'll want to know: \ur{ملک}}
\begin{quote}
\textbf{\ur{ہمارا} \ur{ملک}} — \emph{hamārā mulk} — \textbf{our country.}
\end{quote}

\begin{quote}
\textbf{\ur{اس} \ur{ملک} \ur{میں}} — \emph{is mulk meṅ} — \textbf{in this country.}
\end{quote}

With \textbf{\ur{شہر}} and \textbf{\ur{ملک}} together, a learner can finally answer the two questions every A1 speaking test opens with after a name: \emph{what city}, and \emph{what country}.

\begin{cousinweb}
\textbf{\ur{ملک}} is \textbf{Arabic}, from the three-consonant root \textbf{m-l-k}, \emph{to own, to rule}. The same root gives \textbf{malik}, a king, and the surname \emph{Malik} carried by millions of people from Morocco to Malaysia.

Now put it beside the last lesson. \textbf{\ur{شہر}} is \textbf{Persian} for the seat of a realm; \textbf{\ur{ملک}} is \textbf{Arabic} for the thing a king owns; and \textbf{\ur{کمرہ}} is \textbf{Portuguese} for a vaulted room. Three words for three sizes of place, from three languages, and the reader has met all three in two chapters. That is what Urdu is: a Sanskrit-descended grammar with a Persian, Arabic and European vocabulary poured into it — the same layering that put \emph{khudā} beside \emph{Allāh} in chapter twelve.
\end{cousinweb}

\subsection*{Your turn}
Say these aloud:

\begin{itemize}
\raggedright
  \item \emph{hamārā mulk}
  \item the two scales — \emph{is shahr meṅ}, \emph{is mulk meṅ}
  \item the three languages — \emph{kamrā} Portuguese, \emph{shahr} Persian, \emph{mulk} Arabic
  \item \emph{voh us mulk meṅ hai}
\end{itemize}

\begin{tcolorbox}[breakable,colback=teal!4,colframe=teal!35!black,title={Before you move on}]
What does the root of \textbf{\ur{ملک}} mean? (\textbf{To own, to rule}.) Which king shares it? (\textbf{Malik}.) And which three languages gave this book its three words for a place? (\textbf{Portuguese, Persian and Arabic}.)

Source: \href{https://www.unicode.org/charts/PDF/U0600.pdf}{Unicode Arabic chart}.
\end{tcolorbox}
\section[Lāhaur]{\ur{لاہور} (Lāhaur) — Lahore}

\label{lesson:UR-C27-lahore}



\begin{quote}
Say \textbf{city} and \textbf{country}. Then say \textbf{in this country}.
\end{quote}

\subsection*{You'll want to know: \ur{لاہور}}
\begin{quote}
\textbf{\ur{میں} \ur{لاہور} \ur{میں} \ur{ہوں}} — \emph{maiṅ Lāhaur meṅ hūṅ} — \textbf{I am in Lahore.}
\end{quote}

\begin{quote}
\textbf{\ur{یہ} \ur{لاہور} \ur{ہے}} — \emph{yih Lāhaur hai} — \textbf{this is Lahore.}
\end{quote}

This book has written no place name at all — not a city, not a country, not even the name of the language it teaches. \textbf{\ur{لاہور}} is the first.

Watch the sentence: \textbf{\ur{لاہور} \ur{میں}}, with the postposition after the name and no change to the name itself, because it does not end in \textbf{-\ur{ا}}.

\begin{cousinweb}
Lahore is traditionally \textbf{Lavapuri}, \textquotedblleft{}the city of Lava\textquotedblright{} — Lava being a son of Rama in the Ramayana. So the great city of Urdu letters, where Iqbal wrote and where the language's printing houses stood, carries a \textbf{Sanskrit name out of a Hindu epic}.

That is the same layering the last two lessons showed in the vocabulary, showing up in the map. A Persian word for city, an Arabic word for country, and a Sanskrit name for the place itself.
\end{cousinweb}

\subsection*{Your turn}
Say these aloud:

\begin{itemize}
\raggedright
  \item \emph{yih Lāhaur hai}
  \item \emph{maiṅ Lāhaur meṅ hūṅ}
  \item the three sizes — \emph{is kamre meṅ}, \emph{Lāhaur meṅ}, \emph{is mulk meṅ}
  \item the name behind the name — \emph{Lavapuri}
\end{itemize}

\begin{tcolorbox}[breakable,colback=teal!4,colframe=teal!35!black,title={Before you move on}]
What language is the name \textbf{\ur{لاہور}} from? (\textbf{Sanskrit} — \emph{Lavapuri}.) Does the name change before \textbf{\ur{میں}}? (\textbf{No} — it does not end in \textbf{-\ur{ا}}.) And what is the first place name in this book? (\textbf{This one}.)

Source: \href{https://www.unicode.org/charts/PDF/U0600.pdf}{Unicode Arabic chart}.
\end{tcolorbox}
\section[Lāhaurī]{\ur{لاہوری} (Lāhaurī) — from Lahore}

\label{lesson:UR-C27-lahori}



\begin{quote}
Say \textbf{this is Lahore}. Then say \textbf{heat} and \textbf{cold}, and listen to how both of them end.
\end{quote}

\subsection*{You'll want to know: the ending that makes an origin}
\begin{quote}
\textbf{\ur{میں} \ur{لاہوری} \ur{ہوں}} — \emph{maiṅ Lāhaurī hūṅ} — \textbf{I am from Lahore.}
\end{quote}

Add \textbf{-\ur{ی}} to a place name and you have an adjective meaning \emph{belonging to that place}. That is the whole rule, and this book has already taught you the ending twice without saying so: \emph{garmī} and \emph{sardī} are \emph{garm} and \emph{sard} — \emph{hot} and \emph{cold} — wearing the same \textbf{-\ur{ی}}. (Those four stay in romanization: \emph{garm} needs gāf and \emph{sard} needs dāl, and the script ladder has reached neither.)

\begin{grammarlens}[title={Grammar Lens: the ending that became a surname}]
This is the Arabic \textbf{nisba}, borrowed into Persian and then into Urdu, and once you can see it you will see it everywhere:

\noindent
\begin{tabularx}{\linewidth}{@{}>{\raggedright\arraybackslash}X>{\raggedright\arraybackslash}X@{}}
\toprule
\textbf{name} & \textbf{means} \\
\midrule
\textbf{Bukhārī} & the man from Bukhara \\
\textbf{Ghazālī} & the man from Ghazala \\
\textbf{Lāhaurī} & the one from Lahore \\
\bottomrule
\end{tabularx}

For most of Islamic history a person's last name \textbf{was} their place. The nisba is a naming system, not merely a suffix, and Urdu inherited it whole.

One warning: unlike \textbf{\ur{کالا}} and \textbf{\ur{ہمارا}}, a nisba adjective ends in \textbf{-\ur{ی}} already and does not change for gender. A man and a woman are both \textbf{\ur{لاہوری}}.
\end{grammarlens}

\subsection*{Your turn}
Say these aloud:

\begin{itemize}
\raggedright
  \item \emph{maiṅ Lāhaurī hūṅ}
  \item the ending you already had — \emph{garm}, \emph{garmī}; \emph{sard}, \emph{sardī}
  \item the surnames — \emph{Bukhārī}, \emph{Ghazālī}, \emph{Lāhaurī}
  \item it does not change for gender — \emph{voh Lāhaurī hai}, either way
\end{itemize}

\begin{tcolorbox}[breakable,colback=teal!4,colframe=teal!35!black,title={Before you move on}]
What does adding \textbf{-\ur{ی}} to a place name do? (\textbf{Makes an adjective meaning \emph{from there}}.) Which two words in this book already carried it? (\emph{Garmī} and \emph{sardī}.) And does it change for gender? (\textbf{No}.)

Source: \href{https://www.unicode.org/charts/PDF/U0600.pdf}{Unicode Arabic chart}.
\end{tcolorbox}
\section[shahr · mulk · Lāhaur · Lāhaurī]{Where you are from}

\label{lesson:UR-R27-where-you-are-from}



\begin{quote}
Three words for three sizes of place arrived in this chapter, from three different languages. Name them, and name the language each came from.
\end{quote}

\subsection*{Your turn}
\begin{itemize}
\raggedright
  \item \emph{Say it:} \emph{shahr}, \emph{mulk}, \emph{Lāhaur}
  \item \emph{Say it:} \emph{maiṅ Lāhaur meṅ hūṅ}, then \emph{maiṅ Lāhaurī hūṅ}
  \item \emph{Say it:} the question this answers — \emph{āp kahāṅ …?}
  \item \emph{Recall:} two chapters back, the ending that means \emph{plural, and something is coming} — \emph{-oṅ}
\end{itemize}

\subsection*{Your turn: the distant band — five chapters, and the whole introduction}
Five chapters ago this book could introduce a learner and could not say where they were, where they lived, what belonged to whom, or how many of anything there were. Put the whole introduction together, cold.

Say these aloud:

\begin{itemize}
\raggedright
  \item greeting and calling — \emph{māṅ! salām}
  \item the name — \emph{merā nām … hai}
  \item where you are — \emph{maiṅ Lāhaur meṅ hūṅ}
  \item where you are from — \emph{maiṅ Lāhaurī hūṅ}
  \item the country — \emph{hamārā mulk}
  \item the room, all five shapes — \emph{kamrā}, \emph{kamre}, \emph{kamre meṅ}, \emph{kamroṅ meṅ}, \emph{kamro}
  \item pointing, near and far — \emph{is shahr meṅ}, \emph{us mulk meṅ}, \emph{yahāṅ}, \emph{vahāṅ}
  \item whose — \emph{is kā}, \emph{un kā}, \emph{hamārā}, and \emph{apnā} when it is your own
  \item to whom — \emph{is ko}, \emph{un ko}, \emph{māṅ ko}
  \item and check it — \emph{āp Lāhaurī haiṅ, hai nā?}
  \item offer something, placed — \emph{chāy yahāṅ hai}
\end{itemize}

\begin{tcolorbox}[breakable,colback=teal!4,colframe=teal!35!black,title={Before you move on}]
What turns a place name into \emph{from there}? (\textbf{The -\ur{ی}}.) Which three languages gave this chapter its three words? (\textbf{Persian, Arabic, Sanskrit}.) What must every noun do before a postposition? (\textbf{Take the oblique}.) And what is the one rule that all five of these chapters have been about? (\textbf{A word changes shape when something follows it}.)

Source: \href{https://www.unicode.org/charts/PDF/U0600.pdf}{Unicode Arabic chart}.
\end{tcolorbox}
